% generator/templates/generator/pdf/header.tex
% Attend que main.tex ait défini :
% \typeDS, \ds, \serie, \titre, \baremeGlobal, \sousTitreUn, \sousTitreDeux,
% \headerLineUn, \headerLineDeux, \headerLineTrois, \classe, \nomPrenom, \isCorrection
% ainsi que \ifHAVELOGO et la variable de template {{ header_logo_path }}

\noindent
\begin{tabular*}{\textwidth}{@{\extracolsep{\fill}} p{0.68\textwidth} p{0.28\textwidth}}
  \normalsize \nomPrenom & \raggedleft \normalsize Classe : \classe\\
\end{tabular*}

\vspace{0.25cm}

\noindent
\begin{tikzpicture}[x=1cm,y=1cm, line width=0.5pt]

  % Dimensions
  \pgfmathsetmacro{\X}{\textwidth/1cm}
  \def\H{2}

  % Cadres
  \draw (0,0) rectangle (\X,\H);     % grand rectangle
  \draw (0,0) rectangle (2,\H);      % pavé gauche

  % Pavé droit (barème) uniquement en DS
  \ifx\typeDS\ds
    \draw (\X-2,0) rectangle (\X,\H);
    \draw (\X-2,0) -- (\X,\H);
  \fi

  % Contenu pavé gauche : logo OU 3 lignes
  \ifHAVELOGO
    \node at (1,1) {\includegraphics[width=2cm,height=2cm,keepaspectratio]{\headerLogoPath}};
  \else
    \node[align=center] at (1,1.35) {\large\bfseries \headerLineUn};
    \node[align=center] at (1,1.00) {\small \headerLineDeux};
    \node[align=center] at (1,0.65) {\small \headerLineTrois};
  \fi

  % Abscisse du bloc central (diffère entre DS et Série)
  \ifx\typeDS\ds
    \pgfmathsetmacro{\XC}{\X/2}
  \else
    \pgfmathsetmacro{\XC}{(\X+2)/2}
  \fi

  % Titre et sous-titres
  \node[align=center] at (\XC,1.45) {\Large \bfseries \titre \ifnum\isCorrection=1 ~|~ Correction\fi};
  \node[align=center] at (\XC,0.95) {\normalsize \sousTitreUn};
  \node[align=center] at (\XC,0.55) {\normalsize \sousTitreDeux};

  % Barème (DS)
  \ifx\typeDS\ds
    \node[align=center] at (\X-0.5,0.7) {\Large \bfseries \baremeGlobal};
  \fi

\end{tikzpicture}
